\begingroup

This section removes the perfect-field restriction from normalization
finiteness by constructing finite coefficient-root covers and an actual
separable field container. The Frobenius endomorphism of an arbitrary
ground field is not assumed finite. This also gives the normalization
input for the arbitrary-field curve discussion and reductions in
\S4.3--4.5 of the source.

\begin{lemma}[An actual finite field containing prescribed coefficient roots]
\label{lem:normalization-finite-coefficient-root-field}
\lean{DeJongAlterations.Algebra.IntegralClosureFiniteTypeRoots.coefficientRootField,DeJongAlterations.Algebra.IntegralClosureFiniteTypeRoots.coefficientRootField_finiteDimensional,DeJongAlterations.Algebra.IntegralClosureFiniteTypeRoots.coefficientRoot_pow}
\leanok
Let $k$ be any field, $s\subset k$ finite, and $q>0$. Inside the actual
algebraic closure $\Omega$ of $k$, choose one $q$th root $r_a$ of each
$a\in s$ and set $E=k(r_a:a\in s)$. This is an actual finite intermediate
field of $\Omega/k$, and $r_a^q=a$ holds in $E$ under its canonical
coefficient inclusion.
\end{lemma}
\begin{proof}
The algebraic closure contains a root of $T^q-a$, since its degree is
positive. Every chosen root is algebraic over $k$. An intermediate field
generated by a finite set of algebraic elements is finite dimensional.
Membership in this generated field is the native inclusion of each
chosen root, so its power identity follows by injectivity of the
inclusion into $\Omega$.
\end{proof}

\endgroup

\begingroup

This section removes the perfect-field restriction from normalization
finiteness by constructing finite coefficient-root covers and an actual
separable field container. The Frobenius endomorphism of an arbitrary
ground field is not assumed finite. This also gives the normalization
input for the arbitrary-field curve discussion and reductions in
\S4.3--4.5 of the source.

\begin{lemma}[The actual finite normal polynomial root cover]
\label{lem:normalization-finite-polynomial-root-cover}
\uses{lem:normalization-finite-coefficient-root-field}
\lean{DeJongAlterations.Algebra.IntegralClosureFiniteTypeRoots.polynomialRootMap,DeJongAlterations.Algebra.IntegralClosureFiniteTypeRoots.polynomialRootMap_injective,DeJongAlterations.Algebra.IntegralClosureFiniteTypeRoots.polynomialRootMap_finite,DeJongAlterations.Algebra.IntegralClosureFiniteTypeRoots.polynomialRootCover_isIntegrallyClosed}
\leanok
Let $E/k$ be any actual finite field extension, $\iota$ a finite variable
type, and $q>0$. The literal ring homomorphism
\[
 k[X_i:i\in\iota]\longrightarrow E[Y_i:i\in\iota],
 \qquad a\longmapsto a,\quad X_i\longmapsto Y_i^q,
\]
is injective and finite. Its target is an integrally closed domain.
\end{lemma}
\begin{proof}
The map is coefficient extension followed by the native polynomial
expansion map. Both are injective. Every coefficient of $E$ is integral
over $k$, and remains integral after extending the base to the source
polynomial ring. Each $Y_i$ is integral because its $q$th power is the
image of $X_i$. Polynomial induction therefore proves integrality of
every element of the target. That target is finite type over $k$, since
$E/k$ and the variable type are finite. Its given map from $k$ factors
through the displayed map, so it is also finite type over the source.
Integral plus finite type proves module finiteness of the actual map.
A polynomial ring over the field $E$ is a unique factorization domain,
hence is integrally closed.
\end{proof}

\endgroup

\begingroup

This section removes the perfect-field restriction from normalization
finiteness by constructing finite coefficient-root covers and an actual
separable field container. The Frobenius endomorphism of an arbitrary
ground field is not assumed finite. This also gives the normalization
input for the arbitrary-field curve discussion and reductions in
\S4.3--4.5 of the source.

\begin{lemma}[Actual roots of a prescribed finite polynomial family]
\label{lem:normalization-finite-polynomial-family-roots}
\uses{lem:normalization-finite-coefficient-root-field,lem:normalization-finite-polynomial-root-cover}
\lean{DeJongAlterations.Algebra.IntegralClosureFiniteTypeRoots.familyPolynomialRoot_frobenius,DeJongAlterations.Algebra.IntegralClosureFiniteTypeRoots.familyPolynomialRoot_pow,DeJongAlterations.Algebra.IntegralClosureFiniteTypeRoots.exists_finite_polynomial_root_cover}
\leanok
Let $k$ have exponential characteristic $p$, let $q=p^n$, and let
$\mathcal P$ be a finite family in $k[X_i:i\in\iota]$. There is an
actual finite intermediate field $E\subset\overline k$ such that the
preceding actual polynomial map sends every $P\in\mathcal P$ to a
$q$th power in $E[Y_i:i\in\iota]$.
\end{lemma}
\begin{proof}
Take the finite set of all coefficients occurring in the finite supports
of the given polynomials and construct the preceding root field. For
each $P=\sum_d a_dX^d$, form the actual finite polynomial
$Q_P=\sum_d r_{a_d}Y^d$. The iterated Frobenius is a ring homomorphism;
it sends $Q_P$ to $\sum_da_dY^{qd}$, exactly the image of $P$ under the
coefficient-and-variable root cover. Thus $Q_P^q$ is that image.
\end{proof}

\endgroup

\begingroup

This section removes the perfect-field restriction from normalization
finiteness by constructing finite coefficient-root covers and an actual
separable field container. The Frobenius endomorphism of an arbitrary
ground field is not assumed finite. This also gives the normalization
input for the arbitrary-field curve discussion and reductions in
\S4.3--4.5 of the source.

\begin{lemma}[The actual finite extension of polynomial fraction fields]
\label{lem:normalization-finite-root-fraction-extension}
\uses{lem:normalization-finite-polynomial-root-cover}
\lean{DeJongAlterations.Algebra.IntegralClosureFiniteTypeRoots.fractionRootMap,DeJongAlterations.Algebra.IntegralClosureFiniteTypeRoots.fractionRootMap_algebraMap,DeJongAlterations.Algebra.IntegralClosureFiniteTypeRoots.fractionRootMap_finite}
\leanok
The actual injective polynomial root-cover map extends canonically to
its fraction fields, and this particular fraction-field homomorphism
is finite. It sends $a/b$ to the quotient of the actual images of
$a$ and $b$.
\end{lemma}
\begin{proof}
Use the native universal property of the fraction field for the
injective polynomial root-cover map followed by inclusion in the
target fraction field. Its coefficient equation is exactly that
universal property's equation. The finite source-to-target polynomial
module localizes to a finite module over the source fraction field;
its scalar structure is the same canonical fraction-field lift.
\end{proof}

\endgroup

\begingroup

This section removes the perfect-field restriction from normalization
finiteness by constructing finite coefficient-root covers and an actual
separable field container. The Frobenius endomorphism of an arbitrary
ground field is not assumed finite. This also gives the normalization
input for the arbitrary-field curve discussion and reductions in
\S4.3--4.5 of the source.

\begin{lemma}[Roots of a prescribed finite rational-function family]
\label{lem:normalization-finite-rational-family-roots}
\uses{lem:normalization-finite-polynomial-family-roots,lem:normalization-finite-root-fraction-extension}
\lean{DeJongAlterations.Algebra.IntegralClosureFiniteTypeRoots.exists_finite_rational_root_cover}
\leanok
For any finite family of rational functions in
$\operatorname{Frac}(k[X_i])$ and any $q=p^n$, there is an actual
finite normal polynomial root cover of the preceding kind for which
every selected rational function becomes an actual $q$th power in
the target polynomial fraction field.
\end{lemma}
\begin{proof}
Choose a numerator and a nonzero denominator for each member of the
finite family. The resulting numerator and denominator family is
finite. Apply the polynomial-root construction simultaneously to
this family. The quotient of the chosen polynomial roots has
$q$th power equal to the image of the original rational function,
by the exact fraction-field lift equation and the power-of-quotient
identity. The target polynomial map is injective and finite, as is
its actual fraction-field lift.
\end{proof}

\endgroup

\begingroup

This section removes the perfect-field restriction from normalization
finiteness by constructing finite coefficient-root covers and an actual
separable field container. The Frobenius endomorphism of an arbitrary
ground field is not assumed finite. This also gives the normalization
input for the arbitrary-field curve discussion and reductions in
\S4.3--4.5 of the source.

\begin{lemma}[A rooted Frobenius equation gives actual separability]
\label{lem:normalization-rooted-frobenius-separable}
\lean{DeJongAlterations.Algebra.IntegralClosureFiniteTypeSeparable.eval₂_pow_frobenius,DeJongAlterations.Algebra.IntegralClosureFiniteTypeSeparable.exists_polynomial_frobenius_root,DeJongAlterations.Algebra.IntegralClosureFiniteTypeSeparable.isSeparable_of_rooted_frobenius_equation}
\leanok
Let $F\to E\to\Omega$ be field maps in exponential characteristic
$p$, let $q=p^n$, and let $P\in F[T]$ be separable. Suppose the actual
polynomial $Q\in E[T]$ has coefficientwise $q$th power equal to the
image of $P$. If the image of $P$ vanishes at $x^q\in\Omega$, then
$x$ is separable over $E$.
\end{lemma}
\begin{proof}
Polynomial separability is preserved and reflected by a field map.
Applying this to the coefficient Frobenius and the actual polynomial
identity shows that $Q$ is separable. Frobenius commutes with polynomial
evaluation, so $Q(x)^q=P(x^q)=0$. A field is reduced, hence $Q(x)=0$.
The minimal polynomial of $x$ divides the actual separable polynomial
$Q$, proving the assertion.
\end{proof}

\endgroup

\begingroup

This section removes the perfect-field restriction from normalization
finiteness by constructing finite coefficient-root covers and an actual
separable field container. The Frobenius endomorphism of an arbitrary
ground field is not assumed finite. This also gives the normalization
input for the arbitrary-field curve discussion and reductions in
\S4.3--4.5 of the source.

\begin{lemma}[The actual finite separable field generated by basis images]
\label{lem:normalization-finite-separable-basis-container}
\lean{DeJongAlterations.Algebra.IntegralClosureFiniteTypeSeparable.basisContainer,DeJongAlterations.Algebra.IntegralClosureFiniteTypeSeparable.mem_basisContainer,DeJongAlterations.Algebra.IntegralClosureFiniteTypeSeparable.basisContainerMap,DeJongAlterations.Algebra.IntegralClosureFiniteTypeSeparable.basisContainer_isSeparable,DeJongAlterations.Algebra.IntegralClosureFiniteTypeSeparable.basisContainer_finiteDimensional}
\leanok
Let $L/F$ have a finite basis $(b_i)$, let $F\to E\to\Omega$ be
a field tower, and let $j:L\to\Omega$ be an $F$-embedding. If each
$j(b_i)$ is separable over $E$, the actual field
$M=E(j(b_i):i)$ is finite and separable over $E$. The original map
$j$ factors as a genuine $F$-embedding $L\to M\to\Omega$.
\end{lemma}
\begin{proof}
The finitely many separable generators are algebraic, so their native
generated intermediate field is finite dimensional and separable.
Every $j(x)$ belongs to that field: apply $j$ to the actual finite
basis expansion of $x$. Its coefficients belong to $F$, hence to $E$,
and all basis images are generators. Restricting $j$ to this exact
field gives the asserted injective algebra homomorphism.
\end{proof}

\endgroup

\begingroup

This section removes the perfect-field restriction from normalization
finiteness by constructing finite coefficient-root covers and an actual
separable field container. The Frobenius endomorphism of an arbitrary
ground field is not assumed finite. This also gives the normalization
input for the arbitrary-field curve discussion and reductions in
\S4.3--4.5 of the source.

\begin{lemma}[Finiteness by an actual normal separable cover]
\label{lem:normalization-finite-normal-cover-descent}
\lean{DeJongAlterations.Algebra.IntegralClosureFiniteTypeSeparable.finite_integralClosure_of_finite_normal_separable_cover}
\leanok
Let $A\to B$ be finite, with $A$ Noetherian and $B$ a Noetherian
normal domain. Put $K=\operatorname{Frac}B$. Let $M/K$ be finite
separable and let $j:L\to M$ be an actual $A$-embedding of fields.
Then the actual integral closure of $A$ in $L$ is finite over $A$.
\end{lemma}
\begin{proof}
The native trace theorem makes $C=\operatorname{integralClosure}(B,M)$
finite over $B$, hence over $A$. If $x$ is integral over $A$, then
$j(x)$ is integral over $A$ and therefore over $B$, so it belongs to
that actual $C$. This gives an injective $A$-linear map from the
original integral closure into $C$. A submodule of a finite module
over a Noetherian ring is finite.
\end{proof}

\endgroup

\begingroup

This section removes the perfect-field restriction from normalization
finiteness by constructing finite coefficient-root covers and an actual
separable field container. The Frobenius endomorphism of an arbitrary
ground field is not assumed finite. This also gives the normalization
input for the arbitrary-field curve discussion and reductions in
\S4.3--4.5 of the source.

\begin{theorem}[Arbitrary-field finiteness over the polynomial base]
\label{thm:normalization-finite-polynomial-arbitrary-field}
\uses{lem:normalization-finite-rational-family-roots,lem:normalization-rooted-frobenius-separable,lem:normalization-finite-separable-basis-container,lem:normalization-finite-normal-cover-descent}
\lean{DeJongAlterations.Algebra.IntegralClosureFiniteTypeSeparable.finite_integralClosure_polynomial}
\leanok
Let $k$ be any field, $A=k[X_1,\ldots,X_d]$, $F=\operatorname{Frac}A$,
and $L/F$ a finite field extension. The actual integral closure of $A$
in $L$ is finite over $A$.
\end{theorem}
\begin{proof}
Take the separable closure $F_s$ of $F$ in $L$. The finite purely
inseparable extension $L/F_s$ has an exponent $n$, so every $x^q$,
$q=p^n$, belongs to $F_s$. Choose a finite $F$-basis $(x_i)$ of $L$.
The finitely many minimal polynomials $P_i$ of $x_i^q$ are separable.
Apply the finite rational-function root cover to all their nonzero
coefficients. This gives a finite normal polynomial cover $B$ of $A$
with fraction field $E$ and actual $q$th roots of these coefficients.
The finite polynomials formed with those roots have coefficient
Frobenius equal to the images of $P_i$.
Embed $L$ and $E$ in an actual algebraic closure of $F$. Their maps
are over $F$. In that common field, the preceding lemma shows that
every image of $x_i$ is separable over $E$. Let $M$ be the actual
intermediate field generated over $E$ by those finitely many images.
It is finite and separable over $E$. The image of every element of
$L$ belongs to $M$, by its actual finite basis expansion, giving an
injective map $L\to M$ over the original $A$.
The trace theorem for the integrally closed Noetherian ring $B$
makes its actual integral closure in $M$ finite over $B$, hence over
$A$. Integral elements of $L$ over $A$ map into that actual closure:
their images are integral over $A$, hence over $B$. This is an
injective $A$-linear embedding. Noetherianity of $A$ makes its source
a finite $A$-module.
\end{proof}

\endgroup

\begingroup

This section removes the perfect-field restriction from normalization
finiteness by constructing finite coefficient-root covers and an actual
separable field container. The Frobenius endomorphism of an arbitrary
ground field is not assumed finite. This also gives the normalization
input for the arbitrary-field curve discussion and reductions in
\S4.3--4.5 of the source.

\begin{theorem}[Arbitrary-field finiteness of normalization]
\label{thm:normalization-finite-arbitrary-field}
\uses{thm:normalization-finite-polynomial-arbitrary-field}
\lean{DeJongAlterations.Algebra.IntegralClosureFiniteType.finite_of_isIntegralClosure_of_finiteType,DeJongAlterations.Algebra.IntegralClosureFiniteType.finite_integralClosure,DeJongAlterations.Algebra.IntegralClosureFiniteType.isNoetherianRing_integralClosure,DeJongAlterations.Algebra.IntegralClosureFiniteType.finiteType_integralClosure}
\leanok
Let $D$ be a finite-type domain over any field $k$, and let $L$ be any
fraction field of $D$. Then $\operatorname{integralClosure}(D,L)$
is a finite $D$-module.
\end{theorem}
\begin{proof}
Noether normalization gives an actual finite injective map from a
polynomial ring $A$ over $k$ to $D$. Its fraction field $F$ embeds
in $L$ with $L/F$ finite. The integral closure of $D$ in $L$ is also
the integral closure of $A$ in $L$: integrality of $D/A$ gives one
direction by transitivity, and enlargement of the base gives the
other. The preceding theorem makes this ring finite over $A$.
Restricting the finite scalar generation along $A\to D$ then makes
it finite over $D$.
\end{proof}

\endgroup
